\section{Exact factorization, retained zeros and the failed Cauchy route}
\label{sec:mobius-split-details}
\subsection{Proof of the regrouping with original coefficients}
The monoid consists of nonzero integral ideals prime to the fixed excluded
set $S$. Every divisor of a good ideal is good, so the identity
$\mu*\zeta_0=\delta$ holds on this same monoid. Expanding convolution
in~\eqref{eq:mobius-split} gives finite divisor sums. Multiply the identity
by $\mu(n)^2\nu(n)\chi_n(u)W(Nn/D)$ and sum. The extra squarefree
indicator does not change $\mu(n)$, since $\mu(n)^3=\mu(n)$.
On its support $n=abc$ forces each factor to be squarefree and any two
to be coprime. Sextic multiplicativity, including zeros, gives
$\chi_n(u)=\chi_a(u)\chi_b(u)\chi_c(u)$; a collision with any factor
makes both sides zero. The fixed character also satisfies
$\nu(n)=\nu(a)\nu(b)\nu(c)$ on these good ideals.

In the small-factor piece group $d=ab$. The remaining variable $c$ is
squarefree and coprime to $d$; $Nd\le U^2$ is necessary, but not sufficient
for a nonzero $\beta_U(d)$: its specified truncated factorizations must
still exist. For each such factorization on squarefree $d$,
$\mu(a)\mu(b)=\mu(d)$, proving the coefficient formula.
In the large-factor piece group $s=bc$. The condition $Nb>U$ implies
$Ns>U$, and summing that divisor gives exactly $g_U(s)$.
Since $Ns>U\ge1$, $s$ is not the unit ideal; its complete M\"obius
divisor sum vanishes, proving the alternative formula for $g_U$.
This regrouping never replaces the two-variable weight by $W(Ns/D)$,
never deletes $(a,s)=1$, and never replaces the resulting coefficients
by a favorable arbitrary family.

One can instead keep the unrestricted convolution identity on nonsquarefree
columns, but then both individual pieces and their local conductor
description change. For $U=1$ and a good prime $p$, the column $p^2$
contributes one to each of $a_U*a_U*\zeta_0$ and
$b_U*b_U*\zeta_0$, even though $\mu(p^2)=0$; these cancel only after
the pieces are recombined. Its squared sextic symbol has a cubic local
component. Thus dropping total squarefreeness is not a harmless
piecewise simplification, although the full unrestricted algebraic
identity still holds.

\subsection{The two conductor scales are unchanged}
For a squarefree $n=as$ with $(a,s)=1$, every prime of $n$ has one
nontrivial sextic local component in the row variable. Hence its primitive
finite conductor is the ideal $as$; $NaNs\asymp D$ on the column support.
On multiplying a column by the conjugate of another, the local exponent
is zero at common primes, $1$ at a prime of $n/h$, and $-1$ at a prime
of $n'/h$. The common prime still imposes its original zero on rows.
Writing $n=h r$, $n'=h t$ therefore gives exactly
\[
 \chi_n(u)\overline{\chi_{n'}(u)}
 =\mathbf1_{(u,h)=1}\chi_r(u)\overline{\chi_t(u)},\qquad
 q=rt,\quad\mathfrak l=hrt.
\]
Here $q$ is the primitive oscillatory conductor and $\mathfrak l$ a
period of the full zero extension. These are ideals, not necessarily
squarefree rational norms. In particular, conjugate primes over one split
rational prime are distinct coprime ideal factors. Full defining moduli
of $\nu$ and primary conventions remain fixed in $S$; this row-conductor
calculation does not assert a new uniform reciprocity theorem in the
column variable. In the coprime bulk $Nq\asymp D^2$ remains large.

\subsection{Why a first generic Cauchy step does not supply the saving}
On one type-II dyad let $Na\asymp A$, $Ns\asymp B$, $AB\asymp D$.
Keep the full profile and coprimality restrictions in
\[
 X_a(u)=\nu(a)\chi_a(u)
  \sum_{s\ {\rm squarefree}\atop Ns\asymp B,\,(a,s)=1}
  g_U(s)\nu(s)\chi_s(u)W(NaNs/D).
\]
The exact elementary Cauchy inequality is
\[
 \sum_{0<Nu\le H}\left|D^{-1/2}\sum_{Na\asymp A}\mu(a)X_a(u)\right|^2
 \le \frac{\sum_{Na\asymp A}|\mu(a)|^2}{D}
       \sum_{Na\asymp A}\sum_{0<Nu\le H}|X_a(u)|^2.
\]
The $a$-mask has been kept inside $X_a$. Even if one granted an optimistic
uniform slice bound $(H+B)\sum_{Ns\asymp B}|g_U(s)|^2D^\varepsilon$,
the ordinary divisor estimate $|g_U(s)|\le\tau(s)$ and the fixed-field
ideal count would only yield the envelope
\[
 A(H+B)D^\varepsilon\asymp(AH+D)D^\varepsilon.
\]
That hypothetical slice estimate is not established here. This is an
insufficiency of the generic upper-bound calculation, not a lower bound
on the signed energy. At balanced $A,B\asymp D^{1/2}$ and $H=D^{3/4}$,
the $AH$ term loses $D^{1/2}$ relative to the target; the $D$ term also
exceeds $H$. A sharper use of the actual $g_U$, cross-factor constraints
or covariance between the original pieces is required. This calculation
does not refute all bilinear, moving-cutoff or dispersion methods.

The fixed-$U=1$ lower bounds in
Proposition~\ref{prop:mobius-unit-cutoff} are different: they are genuine
energy lower bounds for the separate pieces. If $W\ge w_0>0$ on
$[c_1,c_2]\subset(1,2)$, count primes in
$[c_1D,c_2D]$ for $I_1$ and in
$[c_1D/Nq_0,c_2D/Nq_0]$ for $II_1$. Removing $S$ and $q_0$ removes
finitely many primes. The prime-ideal theorem gives positive constants
times $D/\log D$ in both cases; the latter contribution has coefficient
two. Because every coefficient of $II_1$ is nonnegative, keeping these
products is valid. This is unlike isolating the prime part of the full
signed M\"obius energy, where uncontrolled cross terms remain.
For the exact requested bound $H(HD)^\varepsilon$ at $H=D^{3/4}$,
the competing power is $D^{3/4+7\varepsilon/4}$, giving the threshold
$\varepsilon<1/7$. Replacing $(HD)^\varepsilon$ by $D^\varepsilon$
reparametrizes the loss and gives the equivalent $1/4$ formulation.
The positivity argument is specific to $U=1$. For a concrete control,
the primary element $7-21\omega$ has norm $637=7\cdot7\cdot13$ and
three distinct ideal prime factors. At $U=7$ its type-I coefficient is
$-1$, its type-II coefficient is $-2$, and its M\"obius value is $-1$.
Thus the separate pieces need not have nonnegative coefficients at another
cutoff; no moving-cutoff obstruction follows from the preceding positivity
proof.

\subsection{Exact finite checks and formal boundaries}
The deterministic checker represents every good ideal of norm at most
$1024$ once by its primary generator, including inert prime ideals and
squarefree products of three distinct prime ideals.
It checks $1535$ coefficient/regrouping cases at five cutoffs, retains
$19680$ encountered original zeros in actual sextic sums, and verifies
$900$ pair-conductor exponent ledgers. It evaluates principal $\nu$ and
the fixed Kummer character $\nu(n)=\chi_n(2)$ on the same good ideals.
The finite diagnostic weight is the rational profile
$((Nn-D)(2D-Nn))^2/D^4$ on $[D,2D]$, zero elsewhere; it is not claimed
to be the smooth analytic weight of the theorem. Every identity is
algebraic for arbitrary finite weights; these samples do not prove
seminorm or scale-supremum estimates.

For example, at $D=128$, $H=16$, principal $\nu$ and $U=1$, the unit-row
normalized sums have the following exact common scaling:
\[
\begin{array}{c|r}
 D^{9/2}M_1&-122721831\\
 D^{9/2}I_1&314726661\\
 D^{9/2}II_1&192004830
\end{array}
\]
Their exact difference is $II_1-I_1=M_1$. The lower-bound theorem uses
the prime-ideal proof, not extrapolation from these small numbers.
An independent no-production-import oracle also checks convolution,
squarefree regrouping and weighted sextic sums.

\path{MobiusSplit.lean} proves eight actual statements: the inverse-ring
identity, natural-number arithmetic truncation/tail, the M\"obius Dirichlet
identity, the unit cutoff, its exact tail convolution and coefficient,
and nonnegativity. The natural-number Dirichlet ring is an actual
arithmetic object; it is not the ideal convolution ring. Ideal arithmetic,
sextic/Hecke character construction, analytic prime counting and the
unproved signed mean square remain prose or named classical inputs.
No new zero-free region, full GRH proof or positive determinant bridge
has been obtained.
